\documentclass{article}
\usepackage{amsmath,amssymb,amsthm}
\usepackage{algorithm}
\usepackage{algpseudocode}
\usepackage{geometry}
\geometry{margin=1in}
\newtheorem{theorem}{Theorem}
\newtheorem{lemma}{Lemma}
\newtheorem{assumption}{Assumption}
\newcommand{\E}{\mathbb{E}}
\newcommand{\R}{\mathbb{R}}
\newcommand{\norm}[1]{\left\|#1\right\|}
\newcommand{\inner}[1]{\left\langle#1\right\rangle}
\begin{document}

\section*{Algorithm Name}
\textbf{Stochastic Gradient Langevin Dynamics with Decreasing Noise (SGLD‑DecVar)}

\section*{Mathematical Formulation}
Let $f:\R^{d}\rightarrow\R$ be the objective function.  
Given a stepsize (learning rate) $\alpha>0$ and a sequence of noise variances
$\{\sigma_{k}^{2}\}_{k\ge 0}$, the iterates $\{x_{k}\}$ are generated by  

\[
\boxed{%
x_{k+1}=x_{k}-\alpha \,\nabla f(x_{k})+\xi_{k}},\qquad k=0,1,2,\dots
\tag{1}
\]

where $\xi_{k}\sim\mathcal N(0,\sigma_{k}^{2}I_{d})$ are independent Gaussian
vectors satisfying  

\[
\E[\xi_{k}]=0,\qquad \E[\|\xi_{k}\|^{2}]=d\,\sigma_{k}^{2}.
\tag{2}
\]

The variance schedule is assumed to be non‑increasing, e.g.  

\[
\sigma_{k}^{2}= \frac{c}{k+1},\qquad c>0.
\tag{3}
\]

\section*{Pseudocode}
\begin{algorithm}[H]
\caption{SGLD‑DecVar}
\begin{algorithmic}[1]
\State \textbf{Input:} stepsize $\alpha>0$, initial point $x_{0}\in\R^{d}$,
          variance constant $c>0$, max iterations $T$.
\For{$k=0$ \textbf{to} $T-1$}
    \State Compute gradient $g_{k}= \nabla f(x_{k})$.
    \State Set variance $\sigma_{k}^{2}=c/(k+1)$.
    \State Sample noise $\xi_{k}\sim\mathcal N(0,\sigma_{k}^{2}I_{d})$.
    \State Update $x_{k+1}=x_{k}-\alpha g_{k}+\xi_{k}$.
\EndFor
\State \textbf{Output:} $\displaystyle \bar x_{T}= \frac1T\sum_{k=0}^{T-1}x_{k}$.
\end{algorithmic}
\end{algorithm}

\section*{Dry Run Example}
Consider $f(w)=(w-3)^{2}$, $w_{0}=0$, stepsize $\alpha=0.1$ and a deterministic
noise $\xi_{k}=0$ (to illustrate the deterministic part).  

\begin{center}
\begin{tabular}{c|c|c|c}
$k$ & $w_{k}$ & $\nabla f(w_{k})=2(w_{k}-3)$ & $w_{k+1}=w_{k}-0.1\nabla f(w_{k})$ \\ \hline
0 & $0$ & $-6$ & $0-0.1(-6)=0.6$ \\
1 & $0.6$ & $2(0.6-3)=-4.8$ & $0.6-0.1(-4.8)=1.08$ \\
2 & $1.08$ & $2(1.08-3)=-3.84$ & $1.08-0.1(-3.84)=1.464$ \\
\end{tabular}
\end{center}
Objective values: $f(0)=9,\;f(0.6)=5.76,\;f(1.08)=3.6864$.

\section*{Assumptions}
\begin{assumption}[Convexity and Smoothness]\label{as:convex}
$f$ is convex and $L$‑smooth, i.e. for all $x,y\in\R^{d}$,
\[
f(y)\le f(x)+\inner{\nabla f(x)}{y-x}+\frac{L}{2}\norm{y-x}^{2}.
\tag{4}
\]
\end{assumption}
\begin{assumption}[Existence of a Minimizer]\label{as:min}
There exists $x^{\star}\in\arg\min f$ and we denote $f^{\star}=f(x^{\star})$.
\end{assumption}
\begin{assumption}[Stepsize]\label{as:step}
The stepsize satisfies $0<\alpha\le \dfrac{1}{2L}$.
\end{assumption}
\begin{assumption}[Noise Schedule]\label{as:noise}
The noise sequence $\{\xi_{k}\}$ is independent, zero‑mean with variances
$\sigma_{k}^{2}=c/(k+1)$ for some $c>0$.
\end{assumption}

\section*{Main Convergence Theorem}
\begin{theorem}\label{thm:conv}
Under Assumptions \ref{as:convex}–\ref{as:noise}, let $\{x_{k}\}$ be generated by
(1) and define the averaged iterate
$\displaystyle \bar x_{T}= \frac{1}{T}\sum_{k=0}^{T-1}x_{k}$. Then for any $T\ge1$,
\[
\E\!\bigl[f(\bar x_{T})-f^{\star}\bigr]\;\le\;
\frac{ \norm{x_{0}-x^{\star}}^{2}+c\bigl(1+\log T\bigr)}{\alpha\,T}.
\tag{5}
\]
Consequently, $\E[f(\bar x_{T})-f^{\star}]=\mathcal O\!\bigl(\tfrac{\log T}{T}\bigr)$.
\end{theorem}

\section*{Theoretical Properties}
\begin{itemize}
\item \textbf{Stability.} The bound (5) holds for any $\alpha\le1/(2L)$,
      guaranteeing that the algorithm does not diverge even with the added
      stochastic perturbation.
\item \textbf{Hyper‑parameter sensitivity.}
      The rate improves linearly with $\alpha$ (larger stepsizes give a
      smaller numerator) but $\alpha$ cannot exceed $1/(2L)$ because the
      contraction factor $1-\alpha L$ would become negative.
\item \textbf{Comparison to SGD.}
      For the same stepsize, SGD with unbiased gradients yields
      $\mathcal O(1/\sqrt{T})$ under bounded variance, whereas the decreasing
      variance schedule of SGLD‑DecVar attains the faster
      $\mathcal O(\log T/T)$ rate for convex smooth objectives.
\end{itemize}

\section*{Discussion}
The theorem shows that the injected Gaussian noise, when decayed as
$c/(k+1)$, does not hinder the convergence of the deterministic gradient
method; instead, it only contributes an additive $\frac{c\log T}{\alpha T}$
term. The result bridges the gap between pure stochastic gradient descent
and Langevin dynamics, providing a principled way to exploit noise for
exploration while preserving optimal convergence for convex problems.

\section*{Preliminaries (Definitions and Lemmas)}
\begin{lemma}[Descent Lemma for $L$‑smooth convex functions]\label{lem:descent}
For any $x\in\R^{d}$,
\[
\|\nabla f(x)\|^{2}\le 2L\bigl(f(x)-f^{\star}\bigr).
\tag{6}
\]
\end{lemma}
\begin{proof}
Convexity gives $f^{\star}\ge f(x)+\inner{\nabla f(x)}{x^{\star}-x}$.  
Smoothness (4) with $y=x^{\star}$ yields
$f^{\star}\le f(x)+\inner{\nabla f(x)}{x^{\star}-x}
+\frac{L}{2}\norm{x^{\star}-x}^{2}$.  
Eliminating the inner product and using $\norm{x^{\star}-x}^{2}\le
\frac{2}{L}\bigl(f(x)-f^{\star}\bigr)$ (a standard consequence of (4))
gives (6). ∎
\end{proof}

\section*{Main Proof (Step‑by‑Step Derivation)}
We prove Theorem \ref{thm:conv}.  
All expectations are taken with respect to the randomness up to iteration
$k$ (including $\xi_{0},\dots,\xi_{k-1}$).

\noindent\textbf{Step 1.} Expand the squared distance after one update:
\[
\begin{aligned}
\norm{x_{k+1}-x^{\star}}^{2}
&=\norm{x_{k}-\alpha\nabla f(x_{k})+\xi_{k}-x^{\star}}^{2}\\
&=\norm{x_{k}-x^{\star}}^{2}
-2\alpha\inner{\nabla f(x_{k})}{x_{k}-x^{\star}}
+\alpha^{2}\norm{\nabla f(x_{k})}^{2}
+2\inner{\xi_{k}}{x_{k}-x^{\star}}
-2\alpha\inner{\nabla f(x_{k})}{\xi_{k}}
+\norm{\xi_{k}}^{2}.
\end{aligned}
\tag{7}
\]

\noindent\textbf{Step 2.} Take conditional expectation $\E[\cdot\,|\,x_{k}]$.
Using $\E[\xi_{k}]=0$ and $\E[\inner{\nabla f(x_{k})}{\xi_{k}}]=0$ we obtain
\[
\E\!\bigl[\norm{x_{k+1}-x^{\star}}^{2}\mid x_{k}\bigr]
= \norm{x_{k}-x^{\star}}^{2}
-2\alpha\inner{\nabla f(x_{k})}{x_{k}-x^{\star}}
+\alpha^{2}\norm{\nabla f(x_{k})}^{2}
+\E\!\bigl[\norm{\xi_{k}}^{2}\bigr].
\tag{8}
\]

\noindent\textbf{Step 3.} Replace $\E[\norm{\xi_{k}}^{2}]$ by the variance
(2):
\[
\E\!\bigl[\norm{\xi_{k}}^{2}\bigr]=d\,\sigma_{k}^{2}.
\tag{9}
\]
For notational simplicity we absorb the dimension $d$ into the constant
$c$ in (3); hence we write $\sigma_{k}^{2}$ directly.

\noindent\textbf{Step 4.} Apply convexity:
\[
\inner{\nabla f(x_{k})}{x_{k}-x^{\star}}
\ge f(x_{k})-f^{\star}.
\tag{10}
\]

\noindent\textbf{Step 5.} Apply Lemma \ref{lem:descent} (6):
\[
\norm{\nabla f(x_{k})}^{2}\le 2L\bigl(f(x_{k})-f^{\star}\bigr).
\tag{11}
\]

\noindent\textbf{Step 6.} Substitute (10) and (11) into (8):
\[
\begin{aligned}
\E\!\bigl[\norm{x_{k+1}-x^{\star}}^{2}\mid x_{k}\bigr]
&\le \norm{x_{k}-x^{\star}}^{2}
-2\alpha\bigl(f(x_{k})-f^{\star}\bigr)
+2\alpha^{2}L\bigl(f(x_{k})-f^{\star}\bigr)
+\sigma_{k}^{2}\\[0.3em]
&= \norm{x_{k}-x^{\star}}^{2}
-2\alpha\bigl(1-\alpha L\bigr)\bigl(f(x_{k})-f^{\star}\bigr)
+\sigma_{k}^{2}.
\end{aligned}
\tag{12}
\]

\noindent\textbf{Step 7.} Use the stepsize restriction (Assumption
\ref{as:step}) $\alpha\le\frac{1}{2L}$, which implies $1-\alpha L\ge\frac12$.
Hence
\[
2\alpha\bigl(1-\alpha L\bigr)\ge\alpha.
\tag{13}
\]

\noindent\textbf{Step 8.} From (12)–(13) we obtain the key recursion:
\[
\E\!\bigl[\norm{x_{k+1}-x^{\star}}^{2}\mid x_{k}\bigr]
\le \norm{x_{k}-x^{\star}}^{2}
-\alpha\bigl(f(x_{k})-f^{\star}\bigr)+\sigma_{k}^{2}.
\tag{14}
\]

\noindent\textbf{Step 9.} Take total expectation on both sides of (14):
\[
\E\!\bigl[\norm{x_{k+1}-x^{\star}}^{2}\bigr]
\le \E\!\bigl[\norm{x_{k}-x^{\star}}^{2}\bigr]
-\alpha\,\E\!\bigl[f(x_{k})-f^{\star}\bigr]
+\sigma_{k}^{2}.
\tag{15}
\]

\noindent\textbf{Step 10.} Rearrange (15) to isolate the suboptimality term:
\[
\alpha\,\E\!\bigl[f(x_{k})-f^{\star}\bigr]
\le \E\!\bigl[\norm{x_{k}-x^{\star}}^{2}\bigr]
-\E\!\bigl[\norm{x_{k+1}-x^{\star}}^{2}\bigr]
+\sigma_{k}^{2}.
\tag{16}
\]

\noindent\textbf{Step 11.} Sum (16) for $k=0,\dots,T-1$:
\[
\alpha\sum_{k=0}^{T-1}\E\!\bigl[f(x_{k})-f^{\star}\bigr]
\le \norm{x_{0}-x^{\star}}^{2}
-\E\!\bigl[\norm{x_{T}-x^{\star}}^{2}\bigr]
+\sum_{k=0}^{T-1}\sigma_{k}^{2}
\le D^{2}+\sum_{k=0}^{T-1}\sigma_{k}^{2},
\tag{17}
\]
where $D^{2}\coloneqq\norm{x_{0}-x^{\star}}^{2}$.

\noindent\textbf{Step 12.} Divide (17) by $\alpha T$ and use convexity of $f$
to replace the average of $f(x_{k})$ by $f(\bar x_{T})$:
\[
\E\!\bigl[f(\bar x_{T})-f^{\star}\bigr]
\le\frac{D^{2}+\sum_{k=0}^{T-1}\sigma_{k}^{2}}{\alpha T}.
\tag{18}
\]

\noindent\textbf{Step 13.} Insert the variance schedule (3):
\[
\sum_{k=0}^{T-1}\sigma_{k}^{2}
= c\sum_{k=0}^{T-1}\frac{1}{k+1}
\le c\bigl(1+\log T\bigr).
\tag{19}
\]

\noindent\textbf{Step 14.} Combine (18) and (19) to obtain the bound
stated in Theorem \ref{thm:conv}:
\[
\E\!\bigl[f(\bar x_{T})-f^{\star}\bigr]
\le\frac{D^{2}+c\bigl(1+\log T\bigr)}{\alpha T}.
\tag{20}
\]
∎

\section*{Rate Derivation (With Constants)}
From (20) we read off the explicit constants:
\[
\boxed{
\text{Rate } = \frac{D^{2}}{\alpha T}
\;+\;
\frac{c}{\alpha}\,\frac{1+\log T}{T}}.
\]
If $c=0$ (no noise) the bound reduces to the classical deterministic
gradient rate $D^{2}/(\alpha T)$.  
With decreasing variance the additional term grows only logarithmically,
hence the overall convergence is $\mathcal O(\log T/T)$.

\section*{Special Cases}
\begin{itemize}
\item \textbf{Constant variance $\sigma_{k}^{2}\equiv\sigma^{2}$.}
      Then $\sum_{k=0}^{T-1}\sigma_{k}^{2}=T\sigma^{2}$ and (20) yields
      $\E[f(\bar x_{T})-f^{\star}]\le
      \frac{D^{2}}{\alpha T}+\frac{\sigma^{2}}{\alpha}$,
      i.e. the error plateaus at $\mathcal O(\sigma^{2})$.
\item \textbf{Strongly convex $f$ with modulus $\mu>0$.}
      Using the stronger inequality
      $\inner{\nabla f(x_{k})}{x_{k}-x^{\star}}\ge
      \mu\norm{x_{k}-x^{\star}}^{2}+f(x_{k})-f^{\star}$,
      the recursion (14) improves to a linear contraction:
      $\E[\norm{x_{k+1}-x^{\star}}^{2}]\le(1-2\alpha\mu) \E[\norm{x_{k}-x^{\star}}^{2}]
      +\sigma_{k}^{2}$,
      leading to an $\mathcal O(1/T)$ rate without the logarithmic factor.
\end{itemize}

\end{document}